\chapter{Introduction}
\label{chpt:intro}

The rapid advancement of Generative Artificial Intelligence (GenAI) is reshaping the
educational landscape, creating new possibilities for personalised learning,
administrative automation, and pedagogical innovation. At the same time, it is
introducing significant uncertainty into existing educational systems, requiring
institutions, policymakers, and educators to reconsider how teaching, learning, and
assessment are designed, governed, and evaluated. As AI tools become more accessible
and increasingly embedded in everyday educational practice, the central challenge is
no longer whether AI will affect education, but how this transformation can be
governed in ways that are pedagogically meaningful, ethically responsible, and
socially equitable.

International organisations have responded to this challenge by developing frameworks
for responsible AI use in education. In particular, \cite{miao2024aiteachers} and
\cite{miao2024aistudents}, in UNESCO's \textit{AI Competency Frameworks for Teachers
and Students}, emphasise that effective AI integration depends not only on the
availability of technological tools, but on teacher capability, student AI literacy,
institutional readiness, ethical awareness, and equity of access. This perspective
reflects a broader shift in the literature, where AI in education is increasingly
understood not as a technical issue, but as a \textit{governance problem} requiring
alignment between policy goals, implementation, and the realities of educational
practice.

However, the literature consistently shows that such alignment remains difficult to
achieve. Across policy and research studies, a recurring challenge concerns the
translation of high-level principles (e.g., transparency, fairness, safety, and
accountability) into practical and context-specific action within educational
settings. In practice, educators are often required to interpret and enact policy in
fast-changing environments without clear procedural guidance, robust training, or
shared institutional norms. This creates a persistent gap between policy aspiration
and educational enactment.

This policy--practice gap is also reflected in up-to-date studies.
\cite{ng2025opportunities} highlights that while AI offers substantial opportunities
for adaptive learning, feedback generation, and administrative efficiency, these
benefits are constrained by concerns regarding teacher preparedness, ethical use,
academic integrity, and unequal access to digital resources. Similar patterns are
evident in the systematic study by \cite{bagdonaite2025artificial}, which shows that
enthusiasm for AI is often accompanied by uncertainty about its pedagogical value,
its implications for student thinking, and the conditions under which its use can be
considered legitimate and educationally sound. Taken together, these findings suggest
that the problem is not simply one of innovation uptake alone, but of implementation
capacity, institutional readiness, and governance coherence.


\begin{center}
  \begin{tcolorbox}[ 
    width=\linewidth, 
    colback=gray!3, 
    colframe=black,
    boxrule=0.6pt, 
    title=Research Problem, 
    colbacktitle=gray!20, 
    coltitle=black,
    fonttitle=\bfseries] 
    Despite the proliferation of AI governance frameworks in education, a persistent
    gap exists between policy design and classroom implementation. This gap is driven
    by limited operational guidance, insufficient teacher preparedness, and unequal
    access to resources, constraining the effective, ethical, and inclusive
    integration of AI in educational practice.
  \end{tcolorbox}
\end{center}

Evidence from policy and public sentiment further underlines this challenge. In
Ireland, for example, recent government guidance by \cite{govie2025ainschools}
acknowledges the need to support schools in managing both the opportunities and the
risks of AI use, especially in relation to assessment, ethics, and adoption. At the
same time, \cite{qqi2025genai} and \cite{ucd2023aiperceptions} suggest a more
nuanced picture: stakeholders recognise AI's potential, but concerns about trust,
regulation, fairness, and practical guidance remain. Comparable patterns also appear
internationally. In the United States, \cite{pew2026aiteens} shows growing
familiarity with AI, especially among younger users, but also concerns around trust,
verification, and the social impact of AI-generated content. Together, these studies
suggest a broader conclusion: successful AI adoption depends not only on policy
design, but also on how such policy is interpreted and enacted by those responsible
for its implementation.


Addressing this problem requires moving beyond descriptive accounts of AI policy 
toward an analysis of how policy is formulated, communicated, interpreted, and 
enacted. The policy--practice gap should be treated not simply as an 
administrative problem, but as a layered governance issue that emerges across 
formal policy documents, institutional responses, and stakeholder discourse. 
This study was therefore guided by the following research question:

\begin{center}
  \begin{tcolorbox}[
      width=\linewidth,
      colback=gray!3,
      colframe=black,
      boxrule=0.6pt,
      title=Research Question,
      colbacktitle=gray!20,
      coltitle=black,
      fonttitle=\bfseries
      ] 
      How can a multi-layered Natural Language Processing (NLP) framework be applied
      to identify, analyse, and explain the policy--practice gap in AI-in-education
      governance across different national contexts?
  \end{tcolorbox}
\end{center}

To answer the research question, this study adopted a computational ``text-as-data''
approach. This choice is grounded in the literature review, which shows that AI
governance, public perceptions, and educational challenges are often examined
separately, despite their close interconnections. The value of a text-as-data
approach lies in its capacity to analyse policy documents and stakeholder discourse
systematically, at a scale that would be difficult to achieve through manual reading
alone. By combining topic modelling, sentiment analysis, and selected elements of
causal inference, the study sought to identify \textit{recurring themes}, detect
\textit{areas of tension or misalignment}, and examine \textit{relationships} between
policy language, stakeholder sentiment, and implementation challenges. In doing so,
it moved beyond a descriptive account of AI policy and developed an integrated
analytical framework for explaining how the policy--practice gap emerges, where it
becomes visible, and how it is sustained across different governance settings.

Drawing on UNESCO's AI Competency Frameworks by \cite{miao2024aiteachers} and
\cite{miao2024aistudents}, which emphasise human-centred, ethical, pedagogical,
technical, and professional dimensions of AI competency, the analysis is organised
around four dimensions: 1) policy and strategy, 2) implementation in practice, 3)
teacher and student competency development, and 4) equity of access. Together, these
dimensions help examine how AI-in-education policy moves from intention toward
practice. They direct attention to the implementation conditions, including
institutional capacity, teacher readiness, learner preparedness, access to resources,
and how policy commitments become actionable in education.

These dimensions are also related to the study of public sentiment. Questions of
trust, fairness, teacher preparedness, access, and ethical use are not only present
in policy frameworks, but also shape how AI in education is perceived by teachers,
students, parents, and wider society. For this reason, the study treated public
sentiment as an important layer of the policy--practice gap. It provides evidence of
how policy priorities are received, questioned, or supported by those affected by AI
adoption in educational settings.

The study adopted a comparative perspective across Ireland, France, Australia, and
the United States. These countries were selected because they offer useful contrasts
in governance structure, policy tradition, and educational response to AI. Ireland is
included as the national context of the study and as a smaller education system
operating within wider European regulatory developments. France provides a
contrasting case because of its more centralised governance model and stronger
emphasis on legal and regulatory structuring. The United States is included because
of its global influence in AI development and its highly decentralised educational
landscape, where implementation is shaped by variation across states, districts, and
institutions. Australia offers an important intermediate case: according to
\cite{australianframeworkgenai2023}, it is a federal system that combines national
coordination with local variation, and it was among the earlier national systems to
develop dedicated guidance on the responsible use of GenAI in schools. Together,
these cases support comparison across centralised, decentralised, and hybrid
governance settings with different implementation capacities.


At the same time, the public sentiment component of the study was approached with a
more focused scope. Because comprehensive and directly comparable public sentiment
data is difficult to collect across all national contexts, the study gave particular
attention to European evidence on public attitudes toward AI and GenAI in education.
This evidence was used as a cautious proxy for broader Western trends, since
\cite{pew2025globalai} shows that European, North American, and Australian contexts
often share concerns around trust, ethical use, regulation, and the social
consequences of AI, while still allowing for national differences in how these
concerns are expressed. To address gaps and imbalance in the available sentiment
evidence, the study also used synthesised sentiment data. This draws on
\cite{imran2022impact} and \cite{bayer2022survey}, which show that synthetic text
generation and data augmentation can support sentiment analysis and text
classification in limited or imbalanced data settings. This supports a more balanced
comparison between policy discourse and public sentiment, provided that systematic
validation is used to assess the faithfulness of the generated data to the original
evidence, including its semantic consistency, thematic alignment, and divergence from
the source data, as discussed by \cite{chim2025evaluating}.

\section*{Study Objectives}

This study was guided by four main objectives:

\begin{enumerate}
    \item Examination of AI-in-education policies and guidance in Ireland, France,
    Australia, and the United States, alongside a review of evidence on public
    sentiment toward AI in education. Given the difficulty of collecting
    comprehensive and directly comparable sentiment data across countries, the study
    gave primary attention to European sentiment evidence, treating it as a cautious
    proxy for broader Western patterns where similar concerns, expectations, and
    debates are likely to emerge.

    \item Assessment of selected NLP and related computational methods, including
    topic modelling, sentiment analysis, elements of causal NLP, and graph-based
    approaches, for linking policy language with stakeholder discourse. This supports
    the aim of examining how policy priorities, public perceptions, and
    implementation challenges relate to one another.

    \item Generation and evaluation of synthetic data to address gaps,
    fragmentation, and limited comparability in available public sentiment datasets.
    Synthetic data was used to support a more balanced relationship between the
    policy and public sentiment corpora, with validation techniques used to assess
    the faithfulness of the synthesised data against the original evidence.

    \item Development of a multi-layered analytical framework for analysing the
    policy--practice gap in AI-in-education governance, with particular attention to
    operational guidance, teacher preparedness, ethical use, and equity of access.
\end{enumerate}


The dissertation was organised as follows. Chapter~\ref{chap:litrev} reviews the
literature on AI governance, public sentiment, and computational approaches to
policy analysis, and identifies the gaps that motivate the study.
Chapter~\ref{chap:method} sets out the research methodology, including the rationale
for the selected analytical methods, experiments, and evaluation criteria.
Chapter~\ref{chap:implresults} describes the implementation of the NLP techniques
and presents the experimental results. Chapter~\ref{chap:analysis} critically
interprets these results in relation to the research question and the wider
literature. Finally, Chapter~\ref{chap:conclusions} summarises the main findings,
discusses the study's limitations, and outlines directions for future research.